\section{A language hard at every sufficiently large length}
\label{diag:section}

The diagonalization follows the classical padded-description method of
Hennie and Stearns \cite[Section~6, Theorem~7]{HennieStearns1966}.
Padding is essential here: each fixed machine must designate itself at
every sufficiently large input length. We give the encoding and simulation
in full to obtain the exact one-tape bounds used in the final transfer.

We use deterministic machines with doubly infinite tapes, finite state sets,
and finite tape alphabets.  The start, accept, and reject states are pairwise
distinct, as are the blank symbol and the two input symbols.  The input is
a binary word written at positions
$0,\ldots,m-1$ of the first tape, all other cells are blank, and every head
starts at position $0$.  A transition moves each head by $-1$, $0$, or $1$.
Acceptance and rejection are designated halting states.  A machine, including
its number of tapes and its alphabet, is fixed independently of the input
length.  We first record the elementary simulation bound used below.

\begin{lemma}[Quadratic simulation]
\label{diag:simulation}
For every fixed deterministic multitape machine $A$ there is a fixed
deterministic one-tape machine $\widehat A$ with the same decisions such that,
if $A$ halts after $T$ steps on an input of length $m$, then $\widehat A$
halts within
\[
                    O_A\bigl((m+T+1)^2\bigr)
\]
steps on that input.
\end{lemma}

\begin{proof}
Suppose that $A$ has $a$ tapes.  One cell of the simulating tape records the
$a$ symbols at the corresponding position of the simulated tapes, together
with a head marker on each track that has its head at that position.  Because
$A$ is fixed, this product alphabet is finite.  Additional marker bits delimit
the represented interval and distinguish old head markers from newly placed
ones.  Initialization takes $O_A(m+1)$ steps.

After $j$ simulated transitions, all nonblank cells and all head positions lie
in an interval of length $O(m+j+1)$.  To perform the next transition, sweep
this interval to read the $a$ marked symbols.  Their tuple and the simulated
state fit in the finite control.  A further sweep performs the indicated
writes and places the new head markers one cell to the left, at the current
cell, or one cell to the right.  Old and new marker bits prevent a marker
from being processed twice.  A further sweep removes the old markers and
promotes the new ones.  Moving to a neighboring cell to place a marker costs
only a bounded number of extra moves.  Extend the represented interval by a
blank cell at either end when necessary.  Thus a fixed number of sweeps,
each taking $O_A(m+j+1)$ steps, implements one transition.  Summing over
$j<T$, including initialization and the final decision, gives
\[
 O_A\bigl((T+1)(m+T+1)\bigr)
       =O_A\bigl((m+T+1)^2\bigr).
\]
The construction neither needs $T$ in advance nor assumes any bound on it.
\end{proof}

We next fix an explicit binary description of one-tape machines.  For a
nonnegative integer $u$, let $\operatorname{un}(u)=1^u0$.  Number the states
$0,\ldots,q-1$ and the tape symbols $0,\ldots,g-1$, where $q,g\geq 3$.
States $0,1,2$ are respectively the start, accept, and reject states.  Tape
symbol $0$ is blank, and symbols $1,2$ encode the input bits $0,1$.
A description begins with $\operatorname{un}(q)\operatorname{un}(g)$.
It then lists, in lexicographic order of the state and scanned symbol, all
$qg$ transition entries.  Each entry consists of the unary codes for the
new state, the written symbol, and a direction code in $\{0,1,2\}$, denoting
left, stationary, and right.  The entries for the halting states are ignored
by the operational semantics.  A description is valid only if it has exactly
these entries, all indices are in range, and it contains no trailing bits.
Every fixed one-tape machine in our model has such a description after
renaming states and symbols; this renaming does not change its running time.

If $b$ is a valid description of length $\ell$, its \emph{designation} is
\[
                        \pi(b)=1^\ell 0b.
\]
A word has at most one designation as a prefix: the initial unary field
specifies exactly how many following bits constitute the description.
Any remaining bits are unrestricted.  Crucially, a machine designated by
a prefix is always simulated on the entire input word, including that
prefix.

\begin{theorem}[A diagonal language]
\label{diag:language}
For every integer $s\geq 1$ there is a binary language $H_s$ with the following
properties.
\begin{enumerate}
\item There is a fixed deterministic one-tape decider $M_s$ and a constant
$C_s$ such that, on every input of length $m\geq 0$, its running time is at
most
\[
                       C_s(m+m^s+2)^{20}.
\]
\item If $D$ is any correct deterministic one-tape decider for $H_s$, then,
for every sufficiently large integer $m$, the worst-case running time of
$D$ on inputs of length $m$ is strictly greater than $m^s$.
\item After their decisions have been recorded, the accepting and rejecting
configurations of $M_s$ may be kept fixed forever.  This gives a stationary
extension of its computation suitable for the local update representation.
\end{enumerate}
\end{theorem}

\begin{proof}
On an input $w$ of length $m$, set $t=m^s$.  If $w$ has no valid designation
as a prefix, put $w\notin H_s$.  Otherwise let $D_b$ be the designated
machine and define
\begin{equation}
\label{diag:definition}
 w\in H_s
 \quad\Longleftrightarrow\quad
 D_b\text{ does not accept the full word }w\text{ within }t\text{ steps}.
\end{equation}
Acceptance within the cutoff includes entering the accepting state on the
last permitted transition.  A rejecting halt never counts as acceptance.

\paragraph{The lower bound.}
Let $D$ be a correct one-tape decider for $H_s$, choose a description $b_D$
of $D$, and write $L_D=|\pi(b_D)|$.  For every $m\geq L_D$ there is an input
of length $m$ beginning with $\pi(b_D)$.  Indeed, every word
$w=\pi(b_D)z$ of that length forces $D$ to run for more than $m^s$ steps.
If $D$ accepted $w$ within the cutoff, Equation~\eqref{diag:definition}
would exclude $w$ from $H_s$.  If $D$ rejected $w$ within the cutoff, then
it would not accept within the cutoff, so the same definition would include
$w$ in $H_s$.  Either possibility contradicts correctness.  Since $D$ is a
decider, its only remaining possibility is to halt after the cutoff.  This
proves the second assertion at every length $m\geq L_D$.

\paragraph{A fixed universal procedure.}
We give the upper bound with substantial slack.  Fix $s$ and write
\[
                          B=m+t+2.
\]
The procedure has a fixed finite number of work tapes and a fixed alphabet;
neither depends on the machine described in the input.  A scan obtains
$m$ in unary and extracts the proposed designation.  Parsing its explicit
table and checking index ranges requires only polynomially many scans of
strings of length $O(m+1)$.  For example, all proposed state and symbol
counts are at most $m$, so even checking every potential entry with unary
counters takes $O((m+1)^5)$ steps.  Malformed inputs are rejected.  When
$m=0$, the input is malformed and this also gives the stated bound.

The cutoff $t$ can be written in unary without any arithmetic oracle.  For
$m\geq 1$, use $s$ nested counters, each ranging over $m$ values, and append
one mark for each counter tuple.  The number of counters is fixed, and a
counter update or reset takes $O_s(m+1)$ time.  This produces exactly $m^s$
marks in $O_s((m+1)(t+1))$ steps.  Copies of these marks supply the loop
budget and the bounds for the tape representation.

For a valid input description, all simulated state and symbol names have
unary length at most $m+1$.  Preallocate records for the tape positions
\[
                         -t-1,\ldots,m+t+1.
\]
There are $m+2t+3=O(B)$ records.  A record contains a symbol field padded
to a fixed width $O(m+1)$ and a head marker; the current simulated state
is stored separately in unary.  Delimiters and padding are symbols of
the fixed simulator alphabet.  The entire configuration list consequently
has length $O(B^2)$.  Initialize positions $0,\ldots,m-1$ with the bits of
the full input $w$, put blanks elsewhere, and mark position $0$.  A head
starting at $0$ cannot leave $[-t,t]$ in $t$ steps, so all simulated accesses
are covered by the preallocated list.

Here is a direct implementation of a simulated transition.  Scan the
configuration list to find its marked record and copy the scanned symbol.
Read the explicit transition table with unary state and symbol counters
until the required entry is found, and copy its new state, symbol, and
direction.  Rescan the configuration list to replace the symbol and move
the marker to the adjacent record, or leave it stationary.  The fixed
field widths ensure that no insertion into the list is necessary.  For a
left move, the simulator may back up one record; for a right move it may
advance one record.  Temporary marker flags distinguish this update from
the next simulated transition.  Finally consume one budget mark and test
whether the new state is accepting or rejecting.

All these operations use scans, copies, and comparisons of explicitly
written strings.  A deliberately inefficient implementation still costs
$O_s(B^5)$ per simulated transition: the configuration list has length
$O(B^2)$, and even a full-list scan for each symbol copied costs $O(B^4)$;
table lookup, comparisons of unary fields, resets, and budget maintenance
fit within one further factor $B$.  Validation, budget construction, and
initialization fit within the same $O_s(B^5)$ bound.  Stop with answer no
as soon as acceptance is detected.  Stop with answer yes on a rejecting
halt or when all $t$ transitions have been simulated without acceptance.
There are at most $t\leq B$ simulated transitions, so this fixed multitape
procedure decides $H_s$ within $O_s(B^6)$ steps.

Apply Lemma~\ref{diag:simulation} to this procedure.  Its one-tape
implementation has time
\[
 O_s\bigl((m+B^6+1)^2\bigr)=O_s(B^{12})
                                      \leq O_s(B^{20}).
\]
Increasing the constant covers every input length and yields the first
assertion.  Denote this fixed one-tape decider by $M_s$.

\paragraph{Stationary halts.}
For subsequent recording of configurations, extend the transition rule of
$M_s$ at either halting state by writing the scanned symbol unchanged,
keeping the head stationary, and keeping the same state.  Reaching the
state still records the decision at the original halting time, but all
later configurations are identical.  The machine has a fixed finite
alphabet and one head, so its configurations can be represented by tape
symbols with an optional head-state marker.  The symbol at a position in
the next configuration is determined by the three symbols at that
position and its two neighbors: the old head may stay there, move in from
a neighbor, or leave after writing.  On valid configurations these
possibilities determine a unique update.  Extending the update arbitrarily
to invalid triples gives the required finite local rule, and the extended
halting configurations are fixed by it.
\end{proof}

\begin{remark}
\label{diag:quantifiers}
The lower-bound threshold in Theorem~\ref{diag:language} may depend on the
decider $D$.  The language and its upper-bound machine depend only on the
fixed integer $s$.  In particular, the theorem rules out running time at
most $m^s$ on all inputs of any length above the threshold associated
with $D$.
\end{remark}
